% !TeX program = tectonic
\documentclass[11pt]{article}
\IfFileExists{alius-guidelines-style.tex}{\usepackage[a4paper,top=27mm,bottom=24mm,left=25mm,right=25mm]{geometry}
\usepackage[T1]{fontenc}
\usepackage[utf8]{inputenc}
\usepackage[scaled=0.96]{helvet}
\renewcommand{\familydefault}{\sfdefault}
\usepackage{array}
\usepackage{enumitem}
\usepackage{fancyhdr}
\usepackage[hidelinks]{hyperref}
\usepackage{microtype}
\usepackage{needspace}
\usepackage{tabularx}
\usepackage{xcolor}
\usepackage{xurl}

\definecolor{aliusgreen}{HTML}{13823A}
\definecolor{aliusgray}{HTML}{7A7A7A}
\definecolor{aliusline}{HTML}{9A9A9A}

\hypersetup{
  colorlinks=true,
  linkcolor=aliusgreen,
  urlcolor=aliusgreen
}

\setlength{\parindent}{0pt}
\setlength{\parskip}{0.72em}
\setlength{\tabcolsep}{0pt}
\raggedbottom
\sloppy
\urlstyle{same}

\newcommand{\ALIUSFooterTitle}{ALIUS Bulletin Guidelines (2022)}

\pagestyle{fancy}
\fancyhf{}
\fancyfoot[L]{\footnotesize\color{aliusgray}\ALIUSFooterTitle}
\fancyfoot[R]{\footnotesize\color{aliusgray}\thepage}
\renewcommand{\headrulewidth}{0pt}
\renewcommand{\footrulewidth}{0.45pt}

\setlist[itemize]{
  leftmargin=1.35em,
  itemsep=0.18em,
  topsep=0.15em,
  parsep=0em
}
\setlist[enumerate]{
  leftmargin=2.05em,
  itemsep=0.42em,
  topsep=0.2em,
  parsep=0em
}
\setlist[description]{
  leftmargin=0em,
  labelsep=0.55em,
  itemsep=0.35em,
  topsep=0.2em,
  font=\normalfont\bfseries
}

\newcommand{\guideDivider}{%
  \par\noindent{\color{aliusline}\rule{\textwidth}{0.45pt}}\par
}

\newcommand{\guideMeta}[2]{%
  {\footnotesize\bfseries\MakeUppercase{#1}}\par
  {\footnotesize #2}\par\vspace{0.65em}
}

\newcommand{\guideHeader}[5]{%
  \thispagestyle{fancy}%
  \vspace*{1mm}
  \begin{center}
    {\Large\bfseries\MakeUppercase{ALIUS Bulletin}\par}
    \vspace{0.25em}
    {\normalsize *\par}
    \vspace{0.25em}
    {\normalsize\bfseries\MakeUppercase{Guidelines For}\par}
    \vspace{0.15em}
    {\LARGE\bfseries\MakeUppercase{#1}\par}
  \end{center}
  \vspace{1.05em}
  \begin{tabularx}{\textwidth}{@{}>{\raggedright\arraybackslash}p{0.55\textwidth}!{\hspace{1.35em}{\color{aliusline}\vrule width 0.5pt}\hspace{1.35em}}>{\raggedright\arraybackslash}X@{}}
    {\large\itshape #2} &
    \guideMeta{Format}{#3}
    \guideMeta{Editorial Route}{#4}
    \guideMeta{Contact}{#5}
  \end{tabularx}
  \vspace{0.75em}
  \guideDivider
  \vspace{0.25em}
}

\newcommand{\guideSection}[1]{%
  \Needspace{6\baselineskip}%
  \par\vspace{1.1em}
  \begin{center}
    {\color{aliusline}\rule{0.14\textwidth}{0.45pt}}\hspace{0.85em}%
    {\large\bfseries #1}%
    \hspace{0.85em}{\color{aliusline}\rule{0.14\textwidth}{0.45pt}}
  \end{center}
  \vspace{-0.25em}
}

\newcommand{\guideSubsection}[1]{%
  \par\vspace{0.65em}{\bfseries #1}\par\vspace{-0.35em}
}

\newcommand{\guideQuestion}[1]{%
  \par\vspace{0.55em}{\color{aliusgreen}\bfseries #1}\par
}

\newenvironment{guideChecklist}
  {\begin{itemize}[label=\textendash]}
  {\end{itemize}}
}{\usepackage[a4paper,top=27mm,bottom=24mm,left=25mm,right=25mm]{geometry}
\usepackage[T1]{fontenc}
\usepackage[utf8]{inputenc}
\usepackage[scaled=0.96]{helvet}
\renewcommand{\familydefault}{\sfdefault}
\usepackage{array}
\usepackage{enumitem}
\usepackage{fancyhdr}
\usepackage[hidelinks]{hyperref}
\usepackage{microtype}
\usepackage{needspace}
\usepackage{tabularx}
\usepackage{xcolor}
\usepackage{xurl}

\definecolor{aliusgreen}{HTML}{13823A}
\definecolor{aliusgray}{HTML}{7A7A7A}
\definecolor{aliusline}{HTML}{9A9A9A}

\hypersetup{
  colorlinks=true,
  linkcolor=aliusgreen,
  urlcolor=aliusgreen
}

\setlength{\parindent}{0pt}
\setlength{\parskip}{0.72em}
\setlength{\tabcolsep}{0pt}
\raggedbottom
\sloppy
\urlstyle{same}

\newcommand{\ALIUSFooterTitle}{ALIUS Bulletin Guidelines (2022)}

\pagestyle{fancy}
\fancyhf{}
\fancyfoot[L]{\footnotesize\color{aliusgray}\ALIUSFooterTitle}
\fancyfoot[R]{\footnotesize\color{aliusgray}\thepage}
\renewcommand{\headrulewidth}{0pt}
\renewcommand{\footrulewidth}{0.45pt}

\setlist[itemize]{
  leftmargin=1.35em,
  itemsep=0.18em,
  topsep=0.15em,
  parsep=0em
}
\setlist[enumerate]{
  leftmargin=2.05em,
  itemsep=0.42em,
  topsep=0.2em,
  parsep=0em
}
\setlist[description]{
  leftmargin=0em,
  labelsep=0.55em,
  itemsep=0.35em,
  topsep=0.2em,
  font=\normalfont\bfseries
}

\newcommand{\guideDivider}{%
  \par\noindent{\color{aliusline}\rule{\textwidth}{0.45pt}}\par
}

\newcommand{\guideMeta}[2]{%
  {\footnotesize\bfseries\MakeUppercase{#1}}\par
  {\footnotesize #2}\par\vspace{0.65em}
}

\newcommand{\guideHeader}[5]{%
  \thispagestyle{fancy}%
  \vspace*{1mm}
  \begin{center}
    {\Large\bfseries\MakeUppercase{ALIUS Bulletin}\par}
    \vspace{0.25em}
    {\normalsize *\par}
    \vspace{0.25em}
    {\normalsize\bfseries\MakeUppercase{Guidelines For}\par}
    \vspace{0.15em}
    {\LARGE\bfseries\MakeUppercase{#1}\par}
  \end{center}
  \vspace{1.05em}
  \begin{tabularx}{\textwidth}{@{}>{\raggedright\arraybackslash}p{0.55\textwidth}!{\hspace{1.35em}{\color{aliusline}\vrule width 0.5pt}\hspace{1.35em}}>{\raggedright\arraybackslash}X@{}}
    {\large\itshape #2} &
    \guideMeta{Format}{#3}
    \guideMeta{Editorial Route}{#4}
    \guideMeta{Contact}{#5}
  \end{tabularx}
  \vspace{0.75em}
  \guideDivider
  \vspace{0.25em}
}

\newcommand{\guideSection}[1]{%
  \Needspace{6\baselineskip}%
  \par\vspace{1.1em}
  \begin{center}
    {\color{aliusline}\rule{0.14\textwidth}{0.45pt}}\hspace{0.85em}%
    {\large\bfseries #1}%
    \hspace{0.85em}{\color{aliusline}\rule{0.14\textwidth}{0.45pt}}
  \end{center}
  \vspace{-0.25em}
}

\newcommand{\guideSubsection}[1]{%
  \par\vspace{0.65em}{\bfseries #1}\par\vspace{-0.35em}
}

\newcommand{\guideQuestion}[1]{%
  \par\vspace{0.55em}{\color{aliusgreen}\bfseries #1}\par
}

\newenvironment{guideChecklist}
  {\begin{itemize}[label=\textendash]}
  {\end{itemize}}
}
\renewcommand{\ALIUSFooterTitle}{ALIUS Bulletin Guidelines for Commentary (2022)}

\begin{document}

\guideHeader
  {Commentary}
  {If you want to edit a Commentary for the ALIUS Bulletin, please read the following guidelines carefully.}
  {Debate commentary or reply; 1,000--2,000 words, plus a 50--100 word abstract.}
  {Published first as an ALIUS Blog article, then integrated into the annual Bulletin.}
  {\href{mailto:aliusresearchgroup@gmail.com}{aliusresearchgroup@gmail.com}}

Debate threads offer a platform to an author, or several authors, to critically discuss a recently published article or book dedicated to consciousness and its altered states.

Each debate thread consists of at least a commentary from the invited author or authors about the target piece. Each debate thread consists of at most: a commentary from the invited author or authors about the target piece; a reply to the commentary from the author or authors of the target piece; a reply to the reply; and a final reply.

Each commentary or reply will be first published as a separate blog article on the ALIUS Blog page, and then subsequently published as part of the ALIUS Bulletin.

\guideSection{Length of the Commentary}

Each Commentary, or reply to the commentary, must include the following:

\begin{guideChecklist}
\item A title.
\item An abstract of 50--100 words.
\item A main text of 1,000--2,000 words.
\item References, with no limit on the number of references. The style used should be the 6th edition of APA. If you use Zotero, this citation style can be downloaded from the \href{https://www.zotero.org/styles?q=id%3Aapa}{Zotero APA style page}.
\item Addresses and affiliations of the authors for correspondence.
\item Font in the submitted file: Times New Roman, size 12.
\end{guideChecklist}

\guideSection{Content of the Commentary}

Comments and replies are not expected to be formulated in a rigorous academic writing style, but rather as an invitation to clarify some positions. Thus, a commentary is not intended to be an attack on the positions defended by the original author. Commentaries can fairly depart from the original approach of the target article in order to offer an alternative viewpoint.

The format also invites the original author to reply and clarify their understanding of the issue, while helping scholars who are non-specialists to better understand misconceptions that appear in the reception of some research. Each comment should be easily understandable to any scholar. For example, an anthropologist reading a comment conducted with a neuroscientist should easily understand most of the discussion, and conversely, a neuroscientist reading an interview conducted with an anthropologist should easily understand most of the discussion.

\guideSection{Modus Operandi and Requirements}

\begin{enumerate}
\item Ask one of the editors of the Bulletin whether the researcher or researchers you want to invite for a commentary fit well with the Bulletin's editorial line.
\item Ask the researcher you want to invite to write a commentary whether they agree in principle to comment. When you do so, let them know about the general concept of the ALIUS Debates, the Bulletin, the aim of ALIUS, the modus operandi, the schedule, and so on.
\item Set up a deadline specifying when the commentary is expected to be received.
\item Once you receive the commentary, invite the original author or authors to write up a reply in the same format and explain to them the general concept of the ALIUS Debates.
\item Once the original author or authors of the target piece have been informed about the publication of the commentary, the commentary will be published on the Blog.
\item If the author or authors of the original target piece have agreed to write up a reply to the commentary, set up a deadline of no more than a month. When the reply is ready, send it to one of the editors for publication on the Blog.
\item Once the reply to the commentary has been received, invite the authors of the commentary to reply to the reply and set up a deadline.
\item Once you have received the reply to the reply, send it to one of the editors and ask the authors of the target piece whether they want to write a final reply. If so, set up a new deadline.
\item When all the back-and-forth commentaries and replies have been collected, you, as editor of the debate thread, should give a title to the whole series. This title should start with the word ``Debating\ldots'' For example: ``Debating the phenomenal vs. access consciousness dichotomy.''
\item Write up short bios of the commentator or commentators and the original author or authors. For length and style, see the bios featured in the first issue of the Bulletin, in the ``Contributors'' section. You may use the official bio provided by the scholar's website, the scholar's university or publisher, and so on.
\item Once the editors have received the commentaries and replies, they will forward them to a native English-speaking proofreader who will check grammar and linguistic correctness. Note that the proofreader's job is not to correct typos. As indicated above, typos must be corrected before sending the file to the editors.
\item The editors will next send you the feedback provided by the proofreaders. You will be expected to integrate the changes suggested by the proofreaders.
\item When the final version is ready, send it to the commentator or commentators and the original author or authors to get their final agreement for publication.
\item Once you have the commentator or commentators' and original author or authors' agreement for publication, send the final version of the debate thread to one of the editors. The comment will then be ready to be published.
\end{enumerate}

\guideSection{Editorial Team}

Daniel Friedman (\href{mailto:danielarifriedman@gmail.com}{danielarifriedman@gmail.com}); Matthieu Koroma (\href{mailto:matthieu.koroma@gmail.com}{matthieu.koroma@gmail.com}); George Fejer (\href{mailto:gyorgyf@live.com}{gyorgyf@live.com}).

\end{document}
